\documentclass[12pt]{article}
\usepackage[T1]{fontenc}
\usepackage[utf8]{inputenc}
\usepackage{lmodern}
\usepackage{textcomp}
\usepackage{enumitem}
\usepackage{geometry}
\geometry{margin=1in}
\begin{document}
\begin{center}
Answer Key - History - Graduate Level Assessment 24 \
\small (Answers are ordered exactly as in the question paper.)
\end{center}

\section*{Multiple Choice}
\begin{enumerate}[label=\arabic*.]
\item \textbf{Answer:} A
\\ \textit{Options:} A) Homo neanderthalensis.; B) Australopithecus afarensis.; C) Australopithecus africanus.; D) Homo naledi.; E) Homo habilis.
\\ \textit{Note:} Homo floresiensis is considered similar or related to Homo neanderthalensis due to shared morphological traits and the possibility of common ancestry within the genus Homo, suggesting a complex evolutionary history among hominins.
\item \textbf{Answer:} E
\\ \textit{Options:} A) music and arts.; B) geography and climate.; C) architecture and construction methods.; D) legal and judicial systems.; E) political and military history.
\\ \textit{Note:} The records left by the Maya primarily focus on political and military history, as seen in their hieroglyphic inscriptions and codices, which document the reigns of rulers, conquests, and significant events, contrasting with the broader economic and administrative concerns often found in Mesopotamian writings.
\item \textbf{Answer:} B
\\ \textit{Options:} A) were primarily animal products, not plant-based.; B) were imported from other regions.; C) were domesticated before groups became sedentary.; D) were domesticated at the same time at communal farms.; E) were mainly wild crops that were not domesticated until the Bronze Age.
\\ \textit{Note:} The correct answer is based on the fact that the Neolithic agricultural revolution in Europe relied heavily on the introduction of domesticated plants and animals from the Near East, as local resources were insufficient to support the diverse agricultural needs of growing populations.
\item \textbf{Answer:} E
\\ \textit{Options:} A) Mousterian technology was focused on the production of scrapers.; B) Aurignacian technology was primarily used for creating stone tools.; C) Aurignacian technology was less advanced than Mousterian technology.; D) Mousterian technology produced more usable blade surface.; E) Aurignacian technology produced more usable blade surface.
\\ \textit{Note:} Aurignacian technology is characterized by its advanced blade production techniques that resulted in longer, more efficient blades with greater usable surface area compared to the predominantly flake-based Mousterian technology.
\item \textbf{Answer:} E
\\ \textit{Options:} A) 500,000 to 100,000 years ago.; B) 800,000 to 200,000 years ago.; C) 200,000 to 50,000 years ago.; D) 700,000 to 40,000 years ago.; E) 600,000 to 30,000 years ago.
\\ \textit{Note:} The correct answer reflects the timeline established by fossil and archaeological evidence indicating that anatomically modern Homo sapiens emerged around 600,000 years ago and persisted until approximately 30,000 years ago, during which they coexisted with other hominin species and developed complex behaviors.
\end{enumerate}

\section*{True / False}
\begin{enumerate}[label=\arabic*.]
\setcounter{enumi}{5}
\item \textbf{Answer:} False
\\ \textit{Options:} A) True; B) False
\\ \textit{Note:} The correct answer is: taswir
\item \textbf{Answer:} False
\\ \textit{Options:} A) True; B) False
\\ \textit{Note:} The correct answer is: World War II
\item \textbf{Answer:} False
\\ \textit{Options:} A) True; B) False
\\ \textit{Note:} The correct answer is: Some of these civilizations
\item \textbf{Answer:} False
\\ \textit{Options:} A) True; B) False
\\ \textit{Note:} The correct answer is: mathematics
\item \textbf{Answer:} False
\\ \textit{Options:} A) True; B) False
\\ \textit{Note:} The correct answer is: French Revolution
\end{enumerate}

\section*{Long Form Response}
\begin{enumerate}[label=\arabic*.]
\setcounter{enumi}{10}
\item \textbf{Answer:} two-phased system
\\ \textit{Note:} Expected answer: two-phased system
\item \textbf{Answer:} le Worm
\\ \textit{Note:} Expected answer: le Worm
\end{enumerate}

\vfill
\noindent\textit{End of Answer Key}
\end{document}